\scanpage{4}
\mspage{\handwritten{7}}
(iv bis) implique $\handwritten{N}=\handwritten{M}_a$, d'où aussitôt que $\text{(iv bis)}\handwritten{\Longrightarrow}\text{(iii ter)}$. Enfin, (iv bis) signifie aussi que $\handwritten{\psi}:X\handwritten{\to}V$ est lisse en $(e^{\handwritten{\bullet}},a)$ et que $\dim X=\dim V$, donc équivaut à (v) en vertu de 1.2. Cela prouve les implications (3). Si $V$ est \deleted{réduit (ou seulement} génériquement réduit\deleted{)}, \deleted{alors}\inserted{comme} (iii bis) implique que la fibré générique de $X\handwritten{\to}V$ est non ramifiée sur le corps rédiduel du point générique de $V$, donc étale sur ce dernier, et que ce corps est aussi l'anneau local de $V$ en son point générique, il s'ensuit que $X\handwritten{\to}V$ est génériquement étale, \deleted{\unsure{il|i.e.|et}} par suite on a \handwritten{(}iii bis)$\handwritten{\Rightarrow}\text{(iii ter)}$. Enfin, si $V$ est normal en $a$, alors $\text{(iv)}\handwritten{\Longrightarrow}\text{(iv bis)}$, car comme $\dim X\handwritten{=}\dim V$ par 1.2. , (iv) implique que $\handwritten{\psi}$ est dominant, et on peut appliquer SGA I 9.5. (ii) qui implique que $\handwritten{\psi}$ est étale et non seulement \inserted{non} ramifié en $(e^{\handwritten{\bullet}},a)$, \deleted{\unsure{ce qui prouv|ce qui prouve}} Cela achève la démonstration de 1.3. Notons d'ailleurs que la condition (ii) implique déjà en vertu de 1.2. que l'inégalité (2) est en fait une égalité. Notons aussi le

\uline{Corolla\deleted{\unsure{o}}ire} 1.4. Supposons que l'on ait $\handwritten{N}=\handwritten{M}_a$\handwritten{,} \inserted{$V$ normal en $a$.} Alors $\phi$ est lisse en $(e,a)$, \deleted{\unsure{(}}donc (1.1. b) $M_a$ est lisse sur $k$ en $e$, et $\deleted{\phi}\handwritten{\psi}$ est lisse en $(e^{\handwritten{\bullet}},a)$.

En effet, on a les inégalités
\[
\dim V+\deleted{\unsure{\dim|\mathrm{r}}}\operatorname{rang}\handwritten{N}\handwritten{\leq}
\dim(G\mathrm{x}W)\handwritten{\leq}\dim_{(e,a)}\phi^{-1}(a)+\dim V
\handwritten{\leq}\dim\handwritten{M}_a+\dim\deleted{\unsure{W|V}}\handwritten{V},
\]
où la première inégalité n'est autre que (2), la deuxième inégalité provient de la \deleted{\unsure{théorie}} théorie de la dimension pour le morphisme \deleted{\unsure{$\phi$}} $\phi:G\mathrm{x}W\handwritten{\to}V$ en le point $(e,a)$, et est une égalité si et seulement si $\phi$ est dominant et la dimension de $\phi^{-1}(a)$ en $(e,a)$ est celle de la fibre générique, ce qui (en vertu de l'hypothèse que $V$ est normal en $a$, et d'un théorème de Chevalley ) signifie aussi que $\phi$ est universellement ouvert en $(e,a)$ ; enfin la troisième inégalité \deleted{est} pro\deleted{\unsure{uv}}\handwritten{vie}nt de l'isomor%
% archive sheet number: 3
